\section{Experimental evaluation}
\label{sec:experiments}
We evaluate \emph{causalKANs} on semi-synthetic benchmarks where ground-truth potential outcomes are available: ACIC 2016 \citep{Dorie2017} and IHDP \citep{hill11bayesian}, settings A and B. Semi-synthetic data allow objective assessment of CATE accuracy despite the fundamental problem of causal inference. We report \itemi the Precision in Estimation of Heterogeneous Effect (PEHE) and \itemii ATE error. Given test set $\mathcal{D}_{\text{test}}$, PEHE is $
% \begin{equation}
% \label{eq:pehe}
% \mathrm{PEHE}\;=\;
% \Bigg(
\sqrt{
\frac{1}{|\mathcal{D}_{\text{test}}|}\sum
% _{\samplecovi_\indexone\in\mathcal{D}_{\text{test}}}
\big(\hcate{\samplecovi_\indexone}-\cate{\samplecovi_\indexone}\big)^2},
% \end{equation}
$
and ATE error is $|\hat{\ate}-\ate|$, where $(\hat{\ate}, {\ate})$ are the estimated and the real ATE, computed as $\ate=\frac{1}{|\mathcal{D}_{\text{test}}|}\sum \cate{\samplecovi}$. All training details follow \cref{sec:causalkans}.

\textbf{Datasets.}
% \emph{IHDP (A/B).} The Infant Health and Development Program is a widely used semi-synthetic benchmark derived from an observational cohort \citep{hill11bayesian}. Outcomes are generated by known functions of the observed covariates, enabling pointwise ground truth. We use the canonical 100 replications, with setting~A (homogeneous effects) and setting~B (nonlinear with heterogeneous effects). The feature dimension is $\covsize{=}25$ and treatment is binary; the sample size is $\samplesize{=}747$ with pronounced class imbalance.  
% \emph{ACIC 2016.} We evaluate on selected realizations from the ACIC’16 challenge, which provide $\samplesize{\approx}4800$ units with $\covsize{=}58$ covariates, binary treatment with covariate-dependent assignment, and nonlinear outcome surfaces. We focus on nonlinear polynomial/exponential regimes (e.g., settings~2/7/26) and use $77$ replications per realization \citep{Dorie2017}. \add{1}{\emph{NSLM}. The national study of learning mindset \citep{yeager2019national} provides an RCT from which \citet{carvalho2019assessing} constructed a semisynthetic dataset for causal inference with featuring over 10,000 observations and 13 covariates. \emph{NEWS.} The news  is a text-based semitynthetic dataset that encodes the count of words in 5000 documents and was used by \citet{johansson16learning} to evaluate causal inference, where the treatment is the device and the outcome is the readers experience. We used the data generating process proposed by \citet{crabbe2022benchmarking}. \emph{TCGA.} We followed the DGP provided by \citet{zhang2023exploring} from the cancer genome atlas dataset \citep{weinstein2013cancer}, which we restricted for having 100 covariates, and binary treatment. More details can be found in \cref{sec:app:datasets}}
\add{1}{\emph{IHDP (A/B).} The IHDP benchmark \citep{hill11bayesian} is constructed from covariates of a real observational study, with synthetic outcomes generated from known functions. We use the standard 100 replications for settings A and B, where A corresponds to a linear outcome surface and B introduces nonlinear and heterogeneous effects.
\emph{ACIC 2016.} The ACIC’16 challenge datasets \citep{Dorie2017} provide covariates from administrative health records and semi-synthetic treatments and outcomes generated through nonlinear mechanisms. We use the nonlinear regimes commonly adopted in prior work.
\emph{NSLM.} The NSLM benchmark uses covariates from the National Study of Learning Mindset \citep{yeager2019national}, with a semi-synthetic DGP introduced by \citet{carvalho2019assessing}.
\emph{NEWS.} The NEWS dataset contains text-derived covariates from a corpus of 5000 documents \citep{johansson16learning}. We adopt the DGP proposed by \citet{crabbe2022benchmarking}.
\emph{TCGA.} The TCGA benchmark uses RNA-seq covariates from the Cancer Genome Atlas \citep{weinstein2013cancer}, combined with the semi-synthetic DGP of \citet{zhang2023exploring}.
Further details for all datasets are provided in \cref{sec:app:datasets}.}
\subsection{Comparison with neural \replace{1}{architectures}{nets}}

\begin{wraptable}{r}{0.5\linewidth}
\vspace{-40pt} % reduce space below

\caption{\add{1}{\textbf{Out-of-sample} }ATE error and PEHE for KAN and MLP across datasets. The baseline (best) according to the Friedman corrected test is \underline{underlined}, and all models not statistically different are \textbf{bolded} ($\significancelevel>=0.05$ in a paired Wilcoxon corrected test). Values are reported as mean\textsubscript{std}.}

\label{tab:kan_vs_mlp_full}
\centering
% \fontsize{7.5pt}{9pt}\selectfont
\scriptsize
\setlength{\tabcolsep}{4pt}
\begin{tabular}{c l cc cc}
% \toprule
& & \multicolumn{2}{c}{KAN} & \multicolumn{2}{c}{MLP} \\
\cmidrule(lr){3-4} \cmidrule(lr){5-6}
Dataset & Model &  ATE err & PEHE & ATE err & PEHE \\
 \hline \\[-1.5ex]
\multirow{4}{*}{IHDP A}
& S-Learner  & \textbf{0.19\textsubscript{0.35}} & \textbf{0.98\textsubscript{1.03}} & \textbf{0.23\textsubscript{0.47}} & 1.04\textsubscript{1.30} \\
& T-Learner  & \textbf{0.13\textsubscript{0.08}} & \underline{\textbf{0.62\textsubscript{0.28}}} & 0.24\textsubscript{0.27} & 0.90\textsubscript{0.51} \\
& TarNet     & \underline{\textbf{0.13\textsubscript{0.10}}} & \textbf{0.64\textsubscript{0.31}} & \textbf{0.17\textsubscript{0.27}} & \textbf{0.70\textsubscript{0.78}} \\
& DragonNet  & \textbf{0.14\textsubscript{0.10}} & \textbf{0.66\textsubscript{0.35}} & \textbf{0.17\textsubscript{0.27}} & \textbf{0.68\textsubscript{0.77}} \\
 \hline \\[-1.5ex]
\multirow{4}{*}{IHDP B}
& S-Learner  & 0.37\textsubscript{0.37} & 3.01\textsubscript{0.63} & 0.32\textsubscript{0.24} & 2.63\textsubscript{0.58} \\
& T-Learner  & 0.34\textsubscript{0.27} & 2.80\textsubscript{0.44} & \textbf{0.25\textsubscript{0.20}} & \textbf{2.06\textsubscript{0.35}} \\
& TarNet     & 0.33\textsubscript{0.24} & 2.68\textsubscript{0.44} & \underline{\textbf{0.23\textsubscript{0.20}}} & \textbf{2.08\textsubscript{0.36}} \\
& DragonNet  & \textbf{0.28\textsubscript{0.22}} & \textbf{2.64\textsubscript{0.44}} & \textbf{0.25\textsubscript{0.21}} & \underline{\textbf{2.00\textsubscript{0.34}}} \\
 \hline \\[-1.5ex]
\multirow{4}{*}{ACIC 2}
& S-Learner  & \underline{\textbf{0.20\textsubscript{0.38}}} & \underline{\textbf{0.20\textsubscript{0.38}}} & 0.38\textsubscript{0.44} & 0.74\textsubscript{0.38} \\
& T-Learner  & 0.56\textsubscript{0.42} & 1.43\textsubscript{0.72} & 1.97\textsubscript{2.13} & 7.05\textsubscript{5.48} \\
& TarNet     & 0.25\textsubscript{0.33} & 0.78\textsubscript{0.44} & 0.36\textsubscript{0.44} & 0.96\textsubscript{0.83} \\
& DragonNet  & \textbf{0.21\textsubscript{0.26}} & 0.75\textsubscript{0.33} & 0.36\textsubscript{0.44} & 0.96\textsubscript{0.82} \\
\hline \\[-1.5ex]
\multirow{4}{*}{ACIC 7}
& S-Learner  & 0.66\textsubscript{0.75} & 4.13\textsubscript{1.56} & 0.75\textsubscript{0.60} & 4.51\textsubscript{1.54} \\
& T-Learner  & \textbf{0.43\textsubscript{0.43}} & \textbf{3.06\textsubscript{1.18}} & 1.37\textsubscript{2.14} & 7.05\textsubscript{6.64} \\
& TarNet     & \underline{\textbf{0.42\textsubscript{0.43}}} & \textbf{3.06\textsubscript{1.18}} & 0.58\textsubscript{0.50} & 4.17\textsubscript{1.47} \\
& DragonNet  & \textbf{0.43\textsubscript{0.43}} & \underline{\textbf{3.06\textsubscript{1.17}}} & 0.59\textsubscript{0.50} & 4.17\textsubscript{1.48} \\
 \hline \\[-1.5ex]
\multirow{4}{*}{ACIC 26}
& S-Learner  & \textbf{0.42\textsubscript{0.45}} & 3.23\textsubscript{1.57} & 0.75\textsubscript{0.60} & 4.51\textsubscript{1.54} \\
& T-Learner  & \textbf{0.36\textsubscript{0.34}} & \textbf{2.80\textsubscript{1.08}} & 1.37\textsubscript{2.14} & 7.05\textsubscript{6.64} \\
& TarNet     & \underline{\textbf{0.35\textsubscript{0.34}}} & \textbf{2.80\textsubscript{1.08}} & 0.58\textsubscript{0.50} & 4.17\textsubscript{1.47} \\
& DragonNet  & \textbf{0.35\textsubscript{0.34}} & \underline{\textbf{2.79\textsubscript{1.08}}} & 0.59\textsubscript{0.50} & 4.17\textsubscript{1.48} \\
 \hline \\[-1.5ex]
\multirow{4}{*}{NSLM}
& S-Learner  & \underline{\textbf{0.048\textsubscript{0.038}}} & \underline{\textbf{0.752\textsubscript{0.008}}} & 0.190\textsubscript{0.032} & 0.753\textsubscript{0.006} \\
& T-Learner  & \textbf{0.050\textsubscript{0.033}} & \underline{\textbf{0.752\textsubscript{0.006}}} & \textbf{0.048\textsubscript{0.033}} & 0.756\textsubscript{0.007} \\
& TarNet     & 0.055\textsubscript{0.034} & \underline{\textbf{0.752\textsubscript{0.006}}} & \textbf{0.049\textsubscript{0.033}} & 0.755\textsubscript{0.007} \\
& DragonNet  & 0.058\textsubscript{0.040} & 0.753\textsubscript{0.007} & \textbf{0.050\textsubscript{0.035}} & 0.756\textsubscript{0.007} \\
 \hline \\[-1.5ex]
\multirow{4}{*}{NEWS}
& S-Learner  & \textbf{2.01\textsubscript{0.68}} & {\textbf{0.141\textsubscript{0.119}}} & 2.73\textsubscript{1.04} & 0.284\textsubscript{0.311} \\
& T-Learner  &  \textbf{2.04\textsubscript{0.53}} & 0.192\textsubscript{0.136} & {2.12\textsubscript{0.54}} & 0.158\textsubscript{0.126} \\
& TarNet & \underline{\textbf{1.97\textsubscript{0.48}}} & 0.157\textsubscript{0.125} & \textbf{2.05\textsubscript{0.44}} & \underline{\textbf{0.120\textsubscript{0.112}}} \\
& DragonNet  & \textbf{2.04\textsubscript{0.72}} & 0.176\textsubscript{0.125} & 2.11\textsubscript{0.44} & \textbf{0.128\textsubscript{0.112}} \\[1ex]
 \hline \\[-1.5ex]
\multirow{4}{*}{\shortstack{TCGA \\ $\times 10^{-2}$}}
& S-Learner  & 2.50\textsubscript{1.50} & 4.64\textsubscript{1.26} & 0.03\textsubscript{0.02} & 2.18\textsubscript{0.27} \\
& T-Learner  & 1.98\textsubscript{6.36} & 4.41\textsubscript{6.27} & \underline{\textbf{0.02\textsubscript{0.01}}} & 2.15\textsubscript{0.07} \\
& TarNet     & 1.19\textsubscript{0.92} & 3.26\textsubscript{0.65} & 0.04\textsubscript{0.01} & \underline{\textbf{1.81\textsubscript{0.03}}} \\
& DragonNet  & 14.95\textsubscript{33.92} & 20.34\textsubscript{33.81} & 0.05\textsubscript{0.04} & 3.74\textsubscript{0.27} \\
\bottomrule
\end{tabular}
% \end{table}
\vspace{-20pt} % reduce space below
\end{wraptable}
\textbf{Baselines.}
We benchmark S-/T-KAN, TARKAN, and DragonKAN against their MLP-based counterparts (S-/T-learner, TARNet, DragonNet). We intentionally restrict baselines to neural counterparts to test our central claim—\emph{comparable accuracy with improved interpretability}—rather than to chase leaderboards. 

We trained causalNNs with the same hyperparameter budget as \ours, modifying depth, number of neurons, activations, regularization and learning rate, and selected the best model based on validation loss (we justify this choice in \cref{sec:app:ablation}).


 \cref{tab:kan_vs_mlp_full} reports PEHE and ATE error per dataset/setting.  We assess statistical significance using the Friedman test with Holm-corrected 
Wilcoxon signed-rank post-hoc comparisons at $\significancelevel = 0.05$ (see \citet{demvsar2006statistical, rainio2024evaluation} for a details on these tests). 
In the tables, the best-performing model (baseline) is \underline{underlined}. 
Models in \textbf{bold} are those for which we cannot reject the null hypothesis 
of equal performance with respect to the baseline ($p \geq 0.05$). A train/val/test split of $80/10/10$ was leveraged, as well as early stopping and Adam optimizer \citep{Kingma2014AdamAM}.


We observe from both \cref{tab:kan_vs_mlp_full} and \cref{fig:pehe_boxplot} that, \replace{1}{across all}{in 7 of the 8} datasets, there exists at least one instance of \our that is the best model or whose performance is statistically indistinguishable from the best-performing neural model,
\clearpage
in the sense that we fail to reject the null hypothesis that the models achieve the same value of the evaluation metric.
\add{1}{On the other hand, in the TCGA dataset, none of the \ours achieve competitive results. However, that does not prevent the use of our pipeline, since one of the steps of it is to compare the metrics against causalNNs. If metrics of \ours are not satisfactory, then the use of KAN based predictors is not recommended by our pipeline (see line 14 of \cref{alg:causalkan}).}

\subsection{Interpretability results}
We show here some examples that highlight the interpretability of \ours, and we will refer to \cref{app:sec:results_visualizations} for more examples and details on these representations. We selected randomly one realization of the shown datasets.

\begin{wrapfigure}{r}{0.5\linewidth}
% \begin{figure}
\centering
\includegraphics[]{figs/all_tkaam_x44.pdf}

\label{sec:interpretability}
\caption{Radar plot and PDP for variable contribution to CATE, using T-KAAM in ACIC-7.}
\label{fig:all_tkaam}
% % \end{figure}
\end{wrapfigure}
\textbf{ACIC-7.} We show results for a single layer \tkan (as it is an additive model, we call it T-KAAM) for ACIC-7, which has been demonstrated very high-performance according to \cref{tab:kan_vs_mlp_full}. This model yields a function of \emph{heterogeneous} CATE, which is a generalized additive model of the covariates (see \cref{sec:app:causalkans_details}). In addition to the closed-form CATE, we provide in \cref{fig:all_tkaam} (extended in \cref{app:sec:results_visualizations_tkaam}) a \emph{Radar plot} \figleft that compares for two individuals, what is the contribution of each variable to the CATE, compared with the average outcome, and a \emph{PDP} \figright that represent the curve that defines the variation of the CATE ($\representationspace$ CATE) with a specific feature. In this case, we can observe how $\covariate_{44}$ increases the CATE w.r.t. the average in individual 10, while it decreases the CATE in individual 11, because the shape of the curve that relates $\covariate_{44}$ and CATE has a maximum near the value of $\covariate_{44}$ of individual 10.

\begin{wrapfigure}{r}{0.4\linewidth}
% \vspace{-40pt}
\centering
\includegraphics[width=0.8\linewidth]{figs/pdp_skaam_treatment.pdf}

\caption{PDP for treatment contribution in $\hpo(\covariates, \treatment)$ estimation, using S-KAAM in IHDP-A.}
\label{fig:skaam_treatment}
\vspace{-20pt}
\end{wrapfigure}
\textbf{IHDP A.} On the other hand, we can estimate \emph{homogeneous} CATEs with a shallow \skan (S-KAAM), which provides closed-form homogeneous treatment effect. In IHDP A, where the known (denoised) ITE is 4 for every individual, we leverage S-KAAM to obtain the value of the causal effect. In \cref{fig:skaam_treatment} we can observe that the predicted formula has an additive linear term relative to the treatment, which correspond to the CATE (as developed in \cref{sec:app:causalkans_details}). The contribution of other variables (not causal) can be consulted in \cref{app:sec:results_visualizations_tkaam}. In this case, the CATE can directly be consulted in the formula: 3.74 for the given example.


% \paragraph{Other datasets.} We omit results for ACIC-26 because they are very similar to those of ACIC-7, while IHDP B and ACIC-2 yield more complex \ours, for which we have not explored visualizations. However, we include in \cref{app:sec:other_datasets} the formulas extracted for one realization of each dataset, which can still be analyzed with function analysis tools.

\subsection{Symbolic CATE recovery}
\label{sec:identifiability}
\add{1}{We have conducted two experiments with known structural causal models, to determine empirically how well \ours capture the equations of the data generating process (DGP). They also provide examples of when the regularization and the simplification steps of the pipeline of \ours help to achieve better approximations to the true data generating functions. Metrics can be consulted in \cref{app:sec:recoverability}}

\add{1}{In \cref{fig:skaam_identifiability}, the function that generates the potential outcomes is additive in the treatment. Therefore, it lies in the class that S-KAAM can model (see \cref{sec:app:causalkans_details}). It can be observed that \itemi the S-KAAM provides a function that is closer to the groundtruth compared with MLP and \itemii the symbolic substitution step approximates the grountruth even better. Note that there is a cosntant value that S-KAAM cannot recover, but does not modify the CATE estimation, since each difference $\hat{\func}(\covariates, \sampletreat_1)-\hat{\func}(\covariates, \sampletreat_1)$ will cancel this constant.}

\add{1}{On the other hand, \cref{fig:tkaam_identifiability} represents the approximation to the CATE with a binary treatment, as the difference between the two potential outcomes, $\{\outcome(1), \outcome(0)\}$. In this case, individual treatment effect is not constant across individuals, and the treatment does not contribute additively to the outcomes. However, the contribution to each variable to the ITE is additive. Therefore, this DGP lies in the class that T-KAAM can model (see \cref{sec:app:causalkans_details}). In the same fashion as in the previous experiment, symbolic T-KAAM is the model that captures better the true function that generates the CATE.}

\begin{figure}[ht]
    \centering
     \begin{subfigure}{0.49\linewidth}
    \centering
        % left: DAG
        \begin{minipage}[c]{0.3\linewidth}
            \centering
            \begin{tikzcd}[
                cells={nodes={main node}}, column sep=small, row sep=small
            ]
                \covariates_1 \arrow[d] \arrow[dr] & \covariates_2  \\
                \treatment \arrow[r] & \outcome\\
            \end{tikzcd}
        \end{minipage}%
        \hfill
        % right: equations
        \begin{minipage}[c]{0.5\linewidth}
            \tiny
            \begin{align*}
                &\covariates_1 \sim \normal(0,1) \\
                &\covariates_2 \sim \normal(0,1) \\
                &\treatment = 1 - 0.1\covariates_1 + \epsilon_{t1} + \epsilon_{t2}, \\
                &\epsilon_{t1} \sim \uniform(-5,5),\quad 
                \epsilon_{t2} \sim \normal(0,1) \\
                &\outcome = 0.5\treatment^{2} + 0.5\treatment + \covariates_1 + \epsilon_\outcome\\
            \end{align*}
        \end{minipage}
    \end{subfigure}
        \begin{subfigure}{0.49\linewidth}
        \centering
        % left: DAG
        \begin{minipage}[c]{0.3\linewidth}
            \centering
            \begin{tikzcd}[
                cells={nodes={main node}}, column sep=small, row sep=small
            ]
                \covariates_1 \arrow[dr] & \covariates_2 \arrow[d]  \\
                \treatment \arrow[r] & \outcome\\
            \end{tikzcd}
        \end{minipage}%
        \hfill
        % right: equations
        \begin{minipage}[c]{0.5\linewidth}
            \tiny
            \begin{align*}
                &\covariates_1 \sim \normal(0,1) \\
                &\covariates_2 \sim \normal(0,1) \\
                &\treatment \sim \normal(0,1) \\
                &\outcome = \treatment \covariates_1 + 0.5 + (1 - \treatment)\covariates_1^{2}
                + 0.5\covariates_2
                + \epsilon_\outcome\\
            \end{align*}
        \end{minipage}
    \end{subfigure}
    \begin{subfigure}{0.49\linewidth}
    \includegraphics[width=\linewidth]{figs/outcome_skaam_first.png}
    \caption{S-KAAM experiment}
    \label{fig:skaam_identifiability}
    \end{subfigure}
    \begin{subfigure}{0.49\linewidth}
    \includegraphics[width=\linewidth]{figs/cate_tkaam.png}
    \caption{T-KAAM experiment}
    \label{fig:tkaam_identifiability}
    \end{subfigure}
    \caption{\add{1}{Synthetic SCMs and the curves recovered by \ours in each step of the pipeline. In blue, the curve recovered by MLP counterparts. Curves green and yellow are overlapped, because the effect of pruning is not appreciable. $\func$ and $\hat{\func}$ represent the groundtruth and the estimated function by Symbolic KAAM, respectively. All curves has been shifted so that the first point concides. \figleft represents directly the equation of the outcome and covariates provided by S-KAAM. \figright represents the equation of the CATE given by T-KAAM, as the difference of two additive formulas.}}
    \label{fig:identifiability_main}
\end{figure}


